% !TEX program = lualatex
% Generated by new_homework.sh — Math 118 Homework 2
\documentclass[11pt]{article}

\usepackage{fontspec}
\usepackage{microtype}
\usepackage{mathtools,amssymb,amsfonts,amsthm,bm}
\usepackage{enumitem}
\usepackage{booktabs,array,xcolor}
\usepackage{geometry,fancyhdr,setspace}
\usepackage{pdfpages}
\usepackage[hidelinks]{hyperref}
\usepackage[nameinlink,noabbrev]{cleveref}

\geometry{margin=1in,headheight=15pt}
\setstretch{1.12}
\setlength{\parindent}{0pt}
\setlength{\parskip}{0.55em plus 0.12em minus 0.08em}
\setlength{\abovedisplayskip}{0.8em plus 0.2em minus 0.1em}
\setlength{\belowdisplayskip}{0.8em plus 0.2em minus 0.1em}
\allowdisplaybreaks[3]
\setlist{leftmargin=*,topsep=0.35em,itemsep=0.15em,parsep=0pt}

\newcommand{\studentname}{Keshav Ramamurthy}
\newcommand{\coursename}{Math 118}
\newcommand{\courselongname}{Fourier Analysis}
\newcommand{\assignment}{Homework 2}
\newcommand{\term}{Fall 2026}

\newcommand{\N}{\mathbb{N}}
\newcommand{\Z}{\mathbb{Z}}
\newcommand{\Q}{\mathbb{Q}}
\newcommand{\R}{\mathbb{R}}
\newcommand{\C}{\mathbb{C}}
\newcommand{\F}{\mathbb{F}}
\newcommand{\K}{\mathbb{K}}
\newcommand{\E}{\mathbb{E}}
\newcommand{\Prob}{\mathbb{P}}
\newcommand{\1}{\mathbf{1}}
\newcommand{\eps}{\varepsilon}
\newcommand{\dd}{\,\mathrm{d}}
\newcommand{\abs}[1]{\left\lvert#1\right\rvert}
\newcommand{\norm}[1]{\left\lVert#1\right\rVert}
\newcommand{\ip}[2]{\left\langle#1,#2\right\rangle}
\newcommand{\set}[1]{\left\{#1\right\}}
\newcommand{\ceil}[1]{\left\lceil#1\right\rceil}
\newcommand{\floor}[1]{\left\lfloor#1\right\rfloor}
\newcommand{\given}{\,\middle\vert\,}
\newcommand{\indep}{\perp\!\!\!\perp}
\newcommand{\wh}[1]{\widehat{#1}}
\newcommand{\deriv}[2]{\frac{\mathrm{d}#1}{\mathrm{d}#2}}
\newcommand{\pderiv}[2]{\frac{\partial#1}{\partial#2}}
\newcommand{\lto}{\longrightarrow}
\newcommand{\st}{\text{ such that }}
\DeclareMathOperator{\Aut}{Aut}
\DeclareMathOperator{\End}{End}
\DeclareMathOperator{\Gal}{Gal}
\DeclareMathOperator{\Hom}{Hom}
\DeclareMathOperator{\Ker}{ker}
\DeclareMathOperator{\im}{im}
\DeclareMathOperator{\rank}{rank}
\DeclareMathOperator{\nullity}{nullity}
\DeclareMathOperator{\Span}{span}
\DeclareMathOperator{\tr}{tr}
\DeclareMathOperator{\diag}{diag}
\DeclareMathOperator{\spec}{spec}
\DeclareMathOperator{\ord}{ord}
\DeclareMathOperator{\supp}{supp}
\DeclareMathOperator{\diam}{diam}
\DeclareMathOperator{\Var}{Var}
\DeclareMathOperator{\Cov}{Cov}
\DeclareMathOperator*{\argmin}{arg\,min}
\DeclareMathOperator*{\argmax}{arg\,max}

\newtheorem{problem}{Problem}
\newtheorem{theorem}{Theorem}
\newtheorem{lemma}{Lemma}
\newtheorem{proposition}{Proposition}
\newtheorem{corollary}{Corollary}
\theoremstyle{definition}
\newtheorem{definition}{Definition}
\newtheorem{example}{Example}
\newtheorem{remark}{Remark}

\newenvironment{solution}{\begin{proof}[Solution]}{\end{proof}}
\newenvironment{answer}{\begin{proof}[Answer]}{\end{proof}}
\renewcommand{\qedsymbol}{$\square$}

\pagestyle{fancy}
\fancyhf{}
\fancyhead[L]{\coursename}
\fancyhead[C]{\assignment\ — my solutions}
\fancyhead[R]{\studentname}
\fancyfoot[C]{\thepage}

\begin{document}

\begin{center}
  {\Large\sffamily\bfseries \coursename\ -- \courselongname}\par\vspace{3pt}
  {\sffamily \assignment\ — my solutions}\par\vspace{3pt}
  {\small\sffamily \term\quad\textbullet\quad \studentname}
\end{center}
\medskip
\hrule
\bigskip

\includepdf[pages=-,pagecommand={\thispagestyle{empty}}]{hw02.pdf}
\clearpage

\section*{Solutions}

\subsubsection*{Problem 1}
(3 points) Let $A$ be an $m\times n$ matrix. Prove that
\[
  \norm{A}_\infty=\max_{1\le i\le m}\sum_{j=1}^n\abs{A_{ij}}.
\]

\begin{solution}
To show that this satisfies the properties of the norm, we go in order.

$\|A\|_\infty\ge0$ and $\|A\|_\infty=0$ implies $A=0$ matrix.

This follows because the maximum sum of the absolute values in each row can only be $0$ if every element of every row is $0$.

\par It is also clear that $\|A\|_\infty\ge0$ because it is a sum of absolute values.

\par We wish to show that $\|cA\|_\infty=|c|\|A\|_\infty$.

\[
\begin{aligned}
\|cA\|_\infty
&= \max_{1\le i\le m}\sum_{j=1}^n |cA_{ij}|\\
&= \max_{1\le i\le m}\sum_{j=1}^n |c||A_{ij}|\\
&= |c|\max_{1\le i\le m}\sum_{j=1}^n |A_{ij}|\\
&= |c|\|A\|_\infty.
\end{aligned}
\]

\par Lastly, we wish to show the triangle inequality. Using the usual triangle inequality,

\[
\begin{aligned}
\|A+B\|_\infty
&= \max_{1\leq i\leq m}\sum_{j=1}^n |A_{ij}+B_{ij}|\\
&\leq \max_{1\leq i\leq m}\sum_{j=1}^n
\left(|A_{ij}|+|B_{ij}|\right)\\
&\leq \max_{1\leq i\leq m}\sum_{j=1}^n |A_{ij}|
+\max_{1\leq i\leq m}\sum_{j=1}^n |B_{ij}|\\
&= \|A\|_\infty+\|B\|_\infty.
\end{aligned}
\]

So all the conditions hold and this is a valid norm.
\end{solution}

\subsubsection*{Problem 2}
(2 points) Suppose $A:V\to W$ and $B:U\to V$ are bounded operators. Show that $(AB)^*=B^*A^*$. If $A$ is invertible, show that $A^*$ is invertible and $(A^*)^{-1}=(A^{-1})^*$.

\begin{solution}
Starting from the adjoint formula, $\langle Tv,w\rangle=\langle v,T^*w\rangle$.

We have $\langle ABv,w\rangle=\langle v,(AB)^*w\rangle$.

\par We can separate $A$ and $B$ using the adjoint formula:

$\langle ABv,w\rangle=\langle Bv,A^*w\rangle=\langle v,B^*A^*w\rangle$.

Thus, $(AB)^*=B^*A^*$.

\par If $A$ is invertible, then $A^{-1}A=I$. Applying the adjoint to both sides gives $A^*(A^{-1})^*=I$.

Similarly, from $AA^{-1}=I$, we get $(A^{-1})^*A^*=I$.

Thus $(A^{-1})^*$ is both a left and right inverse of $A^*$, so $(A^*)^{-1}=(A^{-1})^*$.
\end{solution}

\subsubsection*{Problem 3}
(3 points) (\#12 p.~35. Stop at $x^2$, i.e., drop $x^3$ from the list.) By using the Gram--Schmidt orthogonalization procedure, find an orthonormal basis for the subspace of $L^2[0,1]$ spanned by $1,x,x^2$.

\begin{solution}
We start with $v_1=1$. Its norm is $\sqrt{\int_0^1 1\,dx}=1$, so our first basis vector is $1$.

\par Then, we consider $x$. Using our inner product, $\int_0^1 1x\,dx=\frac{x^2}{2}\big|_0^1=\frac12$, so we subtract off $\frac12$ to get $x-\frac12$.

\par Now we normalize $x-\frac12$:

\[
\sqrt{\int_0^1\left(x-\frac12\right)^2\,dx}
=\frac{1}{2\sqrt3}.
\]

Thus the normalized vector is $2\sqrt3\left(x-\frac12\right)$.

\par Now consider $x^2$. Starting with $1$, we find $\int_0^1 1x^2\,dx=\frac{x^3}{3}\big|_0^1=\frac13$, so we get $x^2-\frac13$.

\par Using the second basis vector,

\[
\int_0^1
2\sqrt3\left(x-\frac12\right)
\left(x^2-\frac13\right)\,dx
=\frac{1}{2\sqrt3}.
\]

Therefore,

$\left(x^2-\frac13\right)-\frac{1}{2\sqrt3}2\sqrt3\left(x-\frac12\right)=x^2-x+\frac16$.

\par Finally, $\int_0^1\left(x^2-x+\frac16\right)^2\,dx=\frac{1}{180}$, so the normalized vector is $6\sqrt5\left(x^2-x+\frac16\right)$.

Thus an orthonormal basis is $1$, $2\sqrt3\left(x-\frac12\right)$, and $6\sqrt5\left(x^2-x+\frac16\right)$.
\end{solution}

\subsubsection*{Problem 4}
(3 points) (\#13 p.~36. Stop at $x^2$ and replace $\cos x$ by $\cos(\pi x)$.) Find the $L^2[0,1]$ projection of the function $\cos(\pi x)$ onto the space spanned by $1,x,x^2$.

\begin{solution}
We want to find the $L^2[0,1]$ projection of $\cos(\pi x)$ onto the space spanned by $1,x,x^2$.

\par Since we already found an orthonormal basis for this space, we can use the projection formula. Let $f(x)=\cos(\pi x)$ and $1$, $2\sqrt3\left(x-\frac12\right)$, and $6\sqrt5\left(x^2-x+\frac16\right)$ be our orthonormal basis.

\par First,

\[
\int_0^1\cos(\pi x)\,dx
=
\left.\frac{\sin(\pi x)}{\pi}\right|_0^1
=
0.
\]

\par Next,

\[
\begin{aligned}
\int_0^1
2\sqrt3\left(x-\frac12\right)\cos(\pi x)\,dx
&=
2\sqrt3\left(
\int_0^1x\cos(\pi x)\,dx
-\frac12\int_0^1\cos(\pi x)\,dx
\right)\\
&=
2\sqrt3\left(-\frac{2}{\pi^2}-0\right)\\
&=
-\frac{4\sqrt3}{\pi^2}.
\end{aligned}
\]

\par Finally,

\[
\begin{aligned}
\int_0^1
6\sqrt5\left(x^2-x+\frac16\right)\cos(\pi x)\,dx
&=
6\sqrt5\left(
\int_0^1x^2\cos(\pi x)\,dx
-\int_0^1x\cos(\pi x)\,dx
+\frac16\int_0^1\cos(\pi x)\,dx
\right)\\
&=
6\sqrt5\left(
-\frac{2}{\pi^2}
+\frac{2}{\pi^2}
+0
\right)\\
&=
0.
\end{aligned}
\]

\par Therefore,

\[
\begin{aligned}
\operatorname{proj}(f)
&=
0(1)
-\frac{4\sqrt3}{\pi^2}
\left[2\sqrt3\left(x-\frac12\right)\right]\\
&=
-\frac{24}{\pi^2}\left(x-\frac12\right).
\end{aligned}
\]
\end{solution}

\subsubsection*{Problem 5}
(3 points) (\#15 p.~36.) Project the function $f(x)=x$ onto the space spanned by
\[
\varphi(x),\ \psi(x),\ \psi(2x),\ \psi(2x-1)\in L^2[0,1],
\]
where
\[
\varphi(x)=
\begin{cases}
1, & 0\le x<1,\\
0, & \text{otherwise},
\end{cases}
\qquad
\psi(x)=
\begin{cases}
1, & 0\le x<1/2,\\
-1, & 1/2\le x<1,\\
0, & \text{otherwise}.
\end{cases}
\]

\begin{solution}
We claim that these functions are all orthogonal. We can see this by computing their pairwise inner products.

\par First,

\[
\langle\varphi,\psi\rangle
=
\int_0^{1/2}(1)(1)\,dx
+
\int_{1/2}^1(1)(-1)\,dx
=
\frac12-\frac12
=
0.
\]

\par Next, $\psi(2x)$ is piecewise $1,-1,0$ when $0\le x<\frac14$, $\frac14\le x<\frac12$, and $\frac12\le x\le1$, respectively. Therefore,

\[
\langle\varphi,\psi(2x)\rangle
=
\int_0^{1/4}(1)(1)\,dx
+
\int_{1/4}^{1/2}(1)(-1)\,dx
+
\int_{1/2}^1(1)(0)\,dx
=
\frac14-\frac14+0
=
0.
\]

\par Similarly, $\psi(2x-1)$ is piecewise $0,1,-1$ when $0\le x<\frac12$, $\frac12\le x<\frac34$, and $\frac34\le x<1$, respectively. Thus,

\[
\langle\varphi,\psi(2x-1)\rangle
=
\int_0^{1/2}(1)(0)\,dx
+
\int_{1/2}^{3/4}(1)(1)\,dx
+
\int_{3/4}^1(1)(-1)\,dx
=
0+\frac14-\frac14
=
0.
\]

\par Now consider $\psi$ and $\psi(2x)$. We get

\[
\langle\psi,\psi(2x)\rangle
=
\int_0^{1/4}(1)(1)\,dx
+
\int_{1/4}^{1/2}(1)(-1)\,dx
+
\int_{1/2}^1(-1)(0)\,dx
=
\frac14-\frac14+0
=
0.
\]

\par Next,

\[
\langle\psi,\psi(2x-1)\rangle
=
\int_0^{1/2}(1)(0)\,dx
+
\int_{1/2}^{3/4}(-1)(1)\,dx
+
\int_{3/4}^1(-1)(-1)\,dx
=
0-\frac14+\frac14
=
0.
\]

\par Finally,

\[
\begin{aligned}
\langle\psi(2x),\psi(2x-1)\rangle
&=
\int_0^{1/4}(1)(0)\,dx
+
\int_{1/4}^{1/2}(-1)(0)\,dx\\
&\qquad+
\int_{1/2}^{3/4}(0)(1)\,dx
+
\int_{3/4}^1(0)(-1)\,dx\\
&=
0.
\end{aligned}
\]

\par Thus they are all orthogonal, and we can add the individual projections of $f(x)=x$ onto each function.

\[
\operatorname{proj}_{\varphi}(f)
=
\frac{\int_0^1x\,dx}{\int_0^1 1\,dx}\varphi
=
\frac12\varphi
\]

\[
\operatorname{proj}_{\psi}(f)
=
\frac{\int_0^{1/2}x\,dx-\int_{1/2}^1x\,dx}
{\int_0^{1/2}1\,dx+\int_{1/2}^1(-1)^2\,dx}\psi
=
-\frac14\psi
\]

\[
\operatorname{proj}_{\psi(2x)}(f)
=
\frac{\int_0^{1/4}x\,dx-\int_{1/4}^{1/2}x\,dx}
{\int_0^{1/4}1\,dx+\int_{1/4}^{1/2}(-1)^2\,dx}\psi(2x)
=
-\frac18\psi(2x)
\]

\[
\operatorname{proj}_{\psi(2x-1)}(f)
=
\frac{\int_{1/2}^{3/4}x\,dx-\int_{3/4}^1x\,dx}
{\int_{1/2}^{3/4}1\,dx+\int_{3/4}^1(-1)^2\,dx}\psi(2x-1)
=
-\frac18\psi(2x-1)
\]

\[
\operatorname{proj}(f)
=
\frac12\varphi
-\frac14\psi
-\frac18\psi(2x)
-\frac18\psi(2x-1).
\]
\end{solution}

\subsubsection*{Problem 6}
(2 points) (\#19 p.~36.) Suppose $K(x,y)$ is a continuous function that vanishes outside a bounded set in $\R^2$. Define $T:L^2(\R)\to L^2(\R)$ by
\[
Tf(x)=\int_{-\infty}^{\infty}K(x,y)f(y)\,dy.
\]

Show that $T^*g(x)=\int_{-\infty}^{\infty}K(y,x)g(y)\,dy$. Note the parallel with the adjoint of a matrix ($A^*_{ij}=A_{ji}$).

\begin{solution}
Starting from the adjoint formula, $\langle Tf,g\rangle=\langle f,T^*g\rangle$.

\[
\begin{aligned}
\langle Tf,g\rangle
&=
\int_{-\infty}^{\infty}Tf(x)g(x)\,dx\\
&=
\int_{-\infty}^{\infty}
\left(\int_{-\infty}^{\infty}K(x,y)f(y)\,dy\right)g(x)\,dx\\
&=
\int_{-\infty}^{\infty}
\int_{-\infty}^{\infty}K(x,y)f(y)g(x)\,dy\,dx\\
&=
\int_{-\infty}^{\infty}
f(y)\left(\int_{-\infty}^{\infty}K(x,y)g(x)\,dx\right)dy.
\end{aligned}
\]

\par Since $x$ is a dummy variable and $y$ is the input variable, we rename the variables to get $T^*g(x)=\int_{-\infty}^{\infty}K(y,x)g(y)\,dy$.
\end{solution}

\subsubsection*{Problem 7}
(2 points) (\#20 p.~36.) Suppose $A:V\to W$ is a linear map between two inner product spaces. Show that $\Ker(A^*)=(\operatorname{Range}(A))^\perp$. Note: Ker stands for Kernel; $\Ker(A^*)$ is the set of all vectors in $W$ which are sent to zero by $A^*$.

\begin{solution}
Let $y\in\ker(A^*)$. Then $A^*y=0$, so for any $x$ using the adjoint formula,

$\langle Ax,y\rangle=\langle x,A^*y\rangle=\langle x,0\rangle=0$.

\par Thus $y$ is orthogonal to every element in $\operatorname{Range}(A)$, so $y\in\operatorname{Range}(A)^\perp$.

\par Therefore, $\ker(A^*)\subseteq\operatorname{Range}(A)^\perp$.

\par Now, let $y\in\operatorname{Range}(A)^\perp$. Then $\langle Ax,y\rangle=0$ for every $x\in V$. Using the adjoint, $\langle x,A^*y\rangle=0$ for every $x\in V$, which implies $A^*y=0$.

Thus $y\in\ker(A^*)$, so $\operatorname{Range}(A)^\perp\subseteq\ker(A^*)$.

\par Therefore, $\ker(A^*)=\operatorname{Range}(A)^\perp$.
\end{solution}

\subsubsection*{Problem 8}
(2 points) (\#21 p.~37. Assume $V$ is finite dimensional --- the theorem is not generally true otherwise.) Prove the following theorem (Fredholm's alternative): Suppose $A:V\to W$ is a linear map between two inner product spaces. Let $b$ be any element in $W$. Then either
\begin{itemize}
  \item $Ax=b$ has a solution $x\in V$; or
  \item there is a vector $w\in W$ with $A^*w=0$ and $\ip{b}{w}_W\ne0$.
\end{itemize}

\begin{solution}
We'll do casework on whether $b$ is in $\operatorname{Range}(A)$.

\par If $b\in\operatorname{Range}(A)$, then by definition there is some $x\in V$ such that $Ax=b$, so the first case holds.

\par Otherwise, $b\notin\operatorname{Range}(A)$. Since $V$ is finite dimensional, $\operatorname{Range}(A)$ is finite dimensional, so $W=\operatorname{Range}(A)\oplus\operatorname{Range}(A)^\perp$.

\par Therefore, since $b\notin\operatorname{Range}(A)$, it has a nonzero orthogonal part. Thus there is some $w\in\operatorname{Range}(A)^\perp$ with $w\neq0$ such that $b=s+w$ for some $s\in\operatorname{Range}(A)$. Therefore,

\[
\langle b,w\rangle_W
=
\langle s+w,w\rangle_W
=
\langle s,w\rangle_W+\langle w,w\rangle_W
=
0+\|w\|^2
>
0.
\]

\par Since $w\in\operatorname{Range}(A)^\perp$, we have $\langle Ax,w\rangle_W=0$ for every $x\in V$. Using the adjoint formula, $\langle x,A^*w\rangle_V=0$ for every $x\in V$, so $A^*w=0$.

\par Therefore, if $b\in\operatorname{Range}(A)$, then $Ax=b$ has a solution. Otherwise, there is a vector $w$ satisfying $A^*w=0$ and $\langle b,w\rangle_W\ne0$.

\par So either $Ax=b$ has a solution, or there is a vector $w$ with $A^*w=0$ and $\langle b,w\rangle_W\ne0$.
\end{solution}

\medskip
\noindent\textit{Page numbers refer to ``A First Course in Wavelets with Fourier Analysis, 2nd Edition'' by Boggess and Narcowich.}

\end{document}
